\subsection{10. Die drei nächsten entscheidenden
Untersuchungen}\label{die-drei-nuxe4chsten-entscheidenden-untersuchungen}

\begin{enumerate}
\def\labelenumi{\arabic{enumi}.}
\tightlist
\item
  \textbf{Den minimalen Ursprungsvertrag auswählen.} Teil B
  klassifiziert die kleinen symmetrieverträglichen Terme und trennt
  echte Erweiterungen von einem bloßen Wechsel der Bosonvariablen. Der
  Compiler muss diese Auswahl und die relativen Parameter begründen.
  Symmetrie allein tut das noch nicht.
\item
  \textbf{Die gemeinsame Schattenkarte am vorhandenen Tensor prüfen.}
  Derselbe Zustand und dieselben Operationen müssen die verschiedenen
  Ausleseansichten tragen. Gesucht sind ergänzende statt nur duplizierte
  Informationen, ein kontrollierter gemeinsamer Kern und Vorhersagen
  ohne nachträgliche Anpassung.
\item
  \textbf{Einen echten Eingriffstransfer und denselben Feldadapter
  konstruieren.} Eine aus der Quelle abgeleitete lokale Operation muss
  später die unbedingte Statistik eines anderen operational bestimmten
  Teils verändern. Selbst zeitabhängige Kreuzkorrelation bei Anfangswert
  null reicht nicht; Teil B enthält ein exaktes Gegenbeispiel. Dazu
  müssen Ladung, Kinetik und relativistische Feldzuordnung passen, bevor
  räumliche Skalierung trägt.
\end{enumerate}

Die korrekte H-Kette mit Restschranke bleibt eine nützliche parallele
Rechenaufgabe. Die zehn bekannten Momente sind feste Gegenprüfungen. Sie
ersetzt aber nicht die jetzt vorrangige Ursprungs- und
Kompositionsfrage.

\subsection{11. Reproduktion und
Veröffentlichung}\label{reproduktion-und-veruxf6ffentlichung}

Das Prüfpaket enthält sechs neue Python-Prüfer, die eingefrorenen
Quellen, die alte und neu wiederholte Kontraktion, Ergebnisdateien,
Quellenmanifest und die vollständige v1.6.6 als historische Basis.
\codeword{replay.py} prüft alle Pins, beide Python-Modi und die
C++-Kontraktion. Die ausgelieferte ZIP wird nach dem Entpacken erneut
geprüft.

Die vollständige Hauptfassung, das kurze Änderungsdokument und die
einfache Erklärung werden gemeinsam versioniert direkt in Documents
abgelegt, jeweils als Markdown und PDF. Die PDF-Ausgaben werden
zusätzlich auf Darstellung, Textabdeckung und Seitenaufbau geprüft.
Frühere Fassungen bleiben erhalten. Es erfolgt keine Änderung der
zentralen TOE-Abnahmemarker, kein Commit, kein Push und keine externe
Web-Veröffentlichung.

Der Statusleitfaden beeinflusst die Trennung zwischen mathematischer
Struktur, angenommenen Ressourcen und nachgewiesener Fähigkeit. Der
Debugging-Leitfaden führte zur ausdrücklichen Sicherung
reell-ganzzahliger Tensorwerte vor der Typumwandlung; keine
Imaginärteile werden still verworfen.

\subsubsection{Schluss in einem Satz}\label{schluss-in-einem-satz}

\textbf{Die einfache Regel ist vorhanden; ihr erster verzweigter
Zustandsbereich und die richtige kleine Rechenkette sind jetzt genauer
gelöst. Was noch fehlt, ist die aus dem Ursprung folgende Auswahl und
Ausführung dieses Prozesses als dieselbe physikalische Welt.}

\begin{center}\rule{0.5\linewidth}{0.5pt}\end{center}

\section{Neue Vogelperspektive: derselbe Prozess in verschiedenen
Ansichten}\label{neue-vogelperspektive-derselbe-prozess-in-verschiedenen-ansichten}

\subsection{B1. Was sich durch die drei weiteren Texte
ändert}\label{b1.-was-sich-durch-die-drei-weiteren-texte-uxe4ndert}

Die nachgereichten Texte werden als Vorschläge geprüft, nicht als
bewiesene Ergebnisse oder neue Arbeitsanweisungen übernommen. Die
stärkste gemeinsame Idee trägt: \textbf{Innere Identität,
Zusammensetzung, Ort, Zeit und Auslese dürfen nicht stillschweigend
dasselbe Label erhalten.} Eine immer genauere Rechnung innerhalb einer
Bank ersetzt diese Verbindungen nicht. Die voranstehenden lokalen
Ergebnisse bleiben gültig; die fundamentale Priorität verschiebt sich
auf die gemeinsame Operations- und Kompositionsregel.

{\def\LTcaptype{none} % do not increment counter
\begin{longtable}[]{@{}
  >{\raggedright\arraybackslash}p{(\linewidth - 2\tabcolsep) * \real{0.5000}}
  >{\raggedright\arraybackslash}p{(\linewidth - 2\tabcolsep) * \real{0.5000}}@{}}
\toprule\noalign{}
\begin{minipage}[b]{\linewidth}\raggedright
Vorschlag
\end{minipage} & \begin{minipage}[b]{\linewidth}\raggedright
Prüfung und Konsequenz
\end{minipage} \\
\midrule\noalign{}
\endhead
\bottomrule\noalign{}
\endlastfoot
Ort in Multiplizitäten statt in den 64 inneren Labels suchen & Sinnvolle
Möglichkeit, aber Multiplizität beschreibt bereits chemische Umwandlung
ohne Raum \\
Der große Kommutant könnte verborgenen Raum enthalten & Hamilton- und
Symmetriekommutant sind verschieden; die Größe allein trägt diesen
Schluss nicht \\
Den ganzen Z4-Ursprung statt nur den ladungserhaltenden Ausschnitt
betrachten & Konkrete zusätzliche Terme existieren am selben W; ihre
Verfügbarkeit ist noch nicht hergeleitet \\
Ladung +3 und Ladung -1 sind modulo vier gleich & Ein zusätzlicher Term
kann sie dynamisch verbinden; die Operatoren sind nicht identisch \\
Zeitabhängige Kreuzantwort zeigt Transport & Im Allgemeinen falsch;
unten steht ein exaktes Gegenbeispiel sogar bei Anfangskorrelation
null \\
Mehrere Schatten könnten einen gemeinsamen Ursprung bestimmen & Möglich
bei gemeinsamem Vertrag und ausreichender Unterscheidbarkeit;
Übereinstimmung allein genügt nicht \\
Holografie könnte die fehlende Verbindung sein & Ein nativer
isometrischer Kodierungsbaustein ist vorhanden; Raumgrenze, Schutz und
gemeinsame Raumzeit fehlen \\
\end{longtable}
}

Die zusätzlich eingefrorenen Texte heißen im Prüfpaket
\codeword{multiplicity\_z4.txt}, \codeword{fundamental\_reduction.txt}
und \codeword{runtime\_synthesis.txt}. Sie werden im Hauptdokument
vollständig angehängt. Die anschließende Schatten- und Holografiefrage
des Nutzers wurde durch eigene Rechnungen weiterverfolgt.

\subsection{B2. Noethers Frage: Welche Unterschiede kann die Dynamik
überhaupt
verändern?}\label{b2.-noethers-frage-welche-unterschiede-kann-die-dynamik-uxfcberhaupt-veruxe4ndern}

Bei ungebrochener innerer Gruppe G zerfällt ein Zustandsraum als

\[\adjustbox{max width=\linewidth}{$\displaystyle 
\mathcal H=\bigoplus_\lambda V_\lambda\otimes M_\lambda,
\qquad H=\bigoplus_\lambda I_{V_\lambda}\otimes h_\lambda,
\qquad \operatorname{End}_G(\mathcal H)
=\bigoplus_\lambda I_{V_\lambda}\otimes\operatorname{End}(M_\lambda).
$}\]

G-invariante Entwicklung darf die Kopien eines Typs verändern, nicht die
inneren Komponenten dieses Typs willkürlich auseinanderziehen. Das macht
Multiplizitäten zu einem möglichen Träger weiterer Dynamik. Es macht sie
noch nicht zu Orten. Schon die drei hellen N=3-Typen enthalten zwei
Kopien: einmal drei Fermionen, einmal einen Fermion plus einen Boson.
Ihre jeweilige Hamiltonmatrix lautet

\[\adjustbox{max width=\linewidth}{$\displaystyle 
h_\lambda=\begin{pmatrix}0&g\sqrt\lambda\\g\sqrt\lambda&\Delta\end{pmatrix},
\qquad \lambda=7,10,12.
$}\]

Am Prüfpunkt g/Delta gleich 1/20 sind die maximalen
Umwandlungswahrscheinlichkeiten 7/107, 1/11 und 3/28. Das ist echte
symmetrieverträgliche Dynamik, zunächst aber eine Veränderung der
Zusammensetzung und kein Ortswechsel.

Wichtig ist die Richtung des Kommutanten: Im N=3-Sektor hat der
Kommutant der vollen inneren Gruppe die Dimension 16. Die viel größere
Zahl 1444233216 betrifft den kleineren Kontrollsatz aus X und Bosonzahl.
Auf dem isolierten irreduziblen 64er-Pol ist der G-Kommutant sogar nur
skalar, während der Kommutant der dort skalaren Energie alle
64-mal-64-Matrizen enthält. Die große Energiedegeneration ist kein
Nachweis vieler Orte.

Für einen G-Singulett-Grundzustand und die ursprünglichen irreduziblen
Fermionoperatoren folgt aus Schurs Lemma

\[\adjustbox{max width=\linewidth}{$\displaystyle 
C_{rs}(t)=\delta_{rs}\,c(t).
$}\]

Feste lineare Umbenennungen liefern lediglich einen Gramfaktor mal
derselben Zeitfunktion. Auch eine perfekte Kenntnis aller lokalen
Spektrallinien würde aus den 64 inneren Komponenten deshalb noch keine
räumliche Anordnung machen. Ein stationärer Grundzustand kann
zeitabhängige Korrelationsfunktionen haben; Nichtstationarität ist keine
notwendige Voraussetzung dafür. Physische Uhrenablesung und Zeitrichtung
sind weitere Fragen.

\subsection{B3. Korrelation ist noch kein Eingriff: ein exakt gelöstes
Gegenbeispiel}\label{b3.-korrelation-ist-noch-kein-eingriff-ein-exakt-geluxf6stes-gegenbeispiel}

Zwei nicht miteinander wechselwirkende Qubits genügen:

\[\adjustbox{max width=\linewidth}{$\displaystyle 
H=\tfrac12(Z_A+Z_B),\qquad
|\Omega\rangle=(|01\rangle+|10\rangle)/\sqrt2,
\qquad A=X_A,\quad B=Y_B.
$}\]

Der Zustand ist stationär. Trotzdem ist

\[\adjustbox{max width=\linewidth}{$\displaystyle 
\langle\Omega|B e^{-itH}A|\Omega\rangle=\sin t.
$}\]

Die Kreuzkorrelation startet bei null und verändert sich mit der Zeit.
Es gibt dennoch keinerlei Signal von A nach B: Alle A-lokalen
Operationen kommutieren mit jedem zeitentwickelten B-lokalen
Observablen. Für jedes spurtreue lokale Quanteninstrument bleibt daher
die unbedingte B-Statistik unverändert. Eine Konditionierung auf ein
Messergebnis in A wäre etwas anderes.

Der belastbare Übertragungstest ist stattdessen

\[\adjustbox{max width=\linewidth}{$\displaystyle 
\delta_{A\to B}(t)=
\operatorname{tr}\!\left[B U_t\mathcal E_A(\rho)U_t^\dagger\right]
-\operatorname{tr}\!\left[B U_t\rho U_t^\dagger\right].
$}\]

A und B müssen zuerst als unabhängig adressierbare Teile definiert sein;
das Instrument muss verfügbar sein, ohne durch seine Definition schon B
anzusteuern. Erfolg bedeutet eine nichtverschwindende unbedingte
Änderung nach dem A-Eingriff. Die lineare Antwort misst entsprechend
einen retardierten Kommutator, nicht bloß eine Zweipunktkorrelation.
Lieb-Robinson-Abschätzungen können anschließend die Ausbreitung in einem
bereits lokal gekoppelten Modell begrenzen; sie erzeugen diese Lokalität
oder drei Raumdimensionen nicht von selbst. Siehe die Primärarbeit von
\href{https://arxiv.org/abs/quant-ph/0603121}{Bravyi, Hastings und
Verstraete}.

Ein weiterer exakter Befund verhindert ein zu starkes Paritäts-No-go:
Der native Graph der 480 Fermionpaare verbindet alle 64 Moden. Eine
Vorzeichenumkehr auf einer Fermionteilmenge, bei unveränderten Bosonen,
erhält jede Kante genau dann, wenn sie trivial oder die globale
Fermionparität ist. Die binäre Inzidenzmatrix hat Rang 63. Zusätzliche
unabhängige Bankparitäten entstehen beim Kopieren von Banken; sie sind
nicht als solche ein Hindernis innerhalb der einen ursprünglichen Bank.
Werden auch Bosonen umgezeichnet, sind andere gemeinsame Gradierungen
möglich.

\subsection{B4. Die einfachste fehlende Auswahl: Warum genau dieser
Hamiltonoperator?}\label{b4.-die-einfachste-fehlende-auswahl-warum-genau-dieser-hamiltonoperator}

Die Z4-Idee besitzt einen präzisen Kern. Ein ganzzahliger, auf jedem
Grad konstanter Ladungslift mit neutralem Grad null müsste bei den
entsprechenden nichtverschwindenden Klammern erfüllen

\[\adjustbox{max width=\linewidth}{$\displaystyle 
2q_1=q_2,\qquad q_1+q_3=0,\qquad 2q_3=q_2.
$}\]

Die Koeffizientenmatrix hat Determinante 4 und Rang drei. Über den
ganzen Zahlen bleibt nur q1=q2=q3=0; modulo vier gilt dagegen die
gewohnte Belegung 1, 2, 3. Das verbietet nicht andere kontinuierliche
Cartanladungen und beweist nicht, dass der aktuelle Operator-Hamiltonian
schon alle E8-Klammern realisiert.

Eine wörtliche Identifikation der Grad-eins-Labels mit primitiven CAR-
Vernichtern scheitert bereits an einem kleinen Beispiel: Die native
W-Spalte des Paares (0,1) ist null, aber \([f_0,f_1]=2f_0f_1\ne0\). Die
innere Lieklammer und die gewöhnliche Operatoralgebra der Fermionfelder
dürfen nicht gleichgesetzt werden. Eine andere, zusammengesetzte
E8-Realisierung wird dadurch nicht ausgeschlossen.

\subsubsection{B4.1 Zwei zusätzliche Terme am unveränderten
Tensor}\label{b4.1-zwei-zusuxe4tzliche-terme-am-unveruxe4nderten-tensor}

Der Bosonträger (10,6) besitzt eine symmetrische reelle invariante
Paarung eta. In der nativen Gewichtsbasis koppelt sie entgegengesetzte
Spin- Vektorgewichte und komplementäre Farbpaare. Es gilt eta Quadrat
gleich eins. Damit sind bereits

\[\adjustbox{max width=\linewidth}{$\displaystyle 
B_+=\tfrac12\sum_{AB}b_A^\dagger\eta_{AB}b_B^\dagger,
\qquad
R_+=\sum_{AB}b_A^\dagger\eta_{AB}P_B^\dagger
$}\]

G-invariant. Bplus erzeugt zwei Bosonen, Rplus einen Boson und zwei
Fermionen. Beide ändern N um vier; beide erhalten N modulo vier. Die
vollständige Lie-Invarianz wurde an allen 45 plus 15 Generatoren
überprüft, nicht nur an den Cartanladungen. Rplus wirkt auf dem leeren
Zustand nicht trivial: Das Normquadrat beträgt exakt 480. Der
quadratische Bosonpaarterm ist sogar einfacher als der vorgeschlagene
zusätzliche kubische Kanal.

Auf dem unveränderten Träger lautet die allgemeine G-invariante,
fermionparitätsgerade, normalgeordnete Hamiltonfamilie vom Grad
höchstens drei

\[\adjustbox{max width=\linewidth}{$\displaystyle 
H_{\rm ext}=c+\varepsilon N_f+\Delta N_b
+\kappa B_++\bar\kappa B_-
+g\sum_A b_A^\dagger P_A+\bar g\sum_A P_A^\dagger b_A
+\lambda R_++\bar\lambda R_-.
$}\]

Die Vollständigkeit gilt nur innerhalb dieser ausdrücklichen
Polynomialklasse: End(F) und End(B) haben jeweils einen Skalar, Sym hoch
zwei von B einen Skalar, und B tritt in Lambda hoch zwei von F genau
einmal auf. Konkret zerfällt letzterer Raum in (10,6), (126,6) und
(120,10). Die SU(4)-Zentrumsregel schließt die verbleibenden
bosonisch-kubischen und gemischten Besetzungs-Boson-Terme aus. Höhere
Polynome, räumliche Kopien und neue Feldträger sind damit nicht
klassifiziert.

Für das native Modell hat das lineare System aller diagonalen
Modenphasen exakt neun unabhängige kontinuierliche Erhaltungen: die acht
inneren Cartans und N. Bei zusätzlichem nichtverschwindendem Rplus-Kanal
oder Bosonpaar-Kanal bleiben exakt acht. Die Ränge sind 115
beziehungsweise 116 auf 124 Moden; explizite ganzzahlige Kerne und
modulare Rangzertifikate beweisen sie. Dies ist keine Behauptung, die
volle nichtabelsche Gruppe schrumpfe auf acht Dimensionen. Sie bleibt
erhalten. Außerdem ist N modulo vier bereits die Wirkung des
SU(4)-Zentrums und keine neu hinzugewonnene unabhängige Symmetrie.

Für reelle Kopplungen und epsilon gleich null enthält die
Fermiongleichung nun zwei verschiedene Beiträge:

\[\adjustbox{max width=\linewidth}{$\displaystyle 
[H_{\rm ext},f_r]=-gD_r-\lambda\widetilde\chi_r^\dagger,
\qquad D_r\sim bf^\dagger,\quad
\widetilde\chi_r^\dagger\sim b^\dagger\eta f^\dagger.
$}\]

Sie haben N-Ladungen -1 und +3. Eine Dynamik kann diese Kanäle bei
bloßer Z4-Erhaltung verbinden, ohne die Operatoren oder ihre
Referenzzustände fälschlich gleichzusetzen. \textbf{Zulässig ist aber
noch nicht aus dem Compiler abgeleitet.} Die Grundzustands- und
Polschranken des unveränderten Modells gelten für diese Erweiterung
nicht ungeprüft weiter.

Auch ein mu-N-Term bleibt G- und Z4-invariant. Bereits der
Zwei-Zustandsblock aus leerem Zustand und einer normierten
Quartetterzeugung im kleinen Modell zeigt

\[\adjustbox{max width=\linewidth}{$\displaystyle 
H+\mu N=\begin{pmatrix}0&\lambda\\\lambda&\Delta+4\mu\end{pmatrix}.
$}\]

Seine niedrigere Eigenenergie hängt von mu ab. Der Übergang zu Z4 allein
wählt also weder ein chemisches Potential noch einen eindeutigen
physikalischen Grundzustand aus. Für Delta größer als Betrag kappa ist
die bosonische Quadratik positiv; die linearen Bosonkopplungen an
beschränkte Fermionoperatoren lassen sich relativ dazu abschätzen. Das
gibt eine untere Energieschranke, aber noch keinen neuen eindeutigen
Grundzustandssatz.

\subsubsection{B4.2 Eine weitere Vereinfachung: neue Dynamik oder andere
Teilchenvariablen?}\label{b4.2-eine-weitere-vereinfachung-neue-dynamik-oder-andere-teilchenvariablen}

Manche scheinbaren Erweiterungen beschreiben nur andere Bosonvariablen.
Auf der reellen, phasengleichen Teilfamilie setze man

\[\adjustbox{max width=\linewidth}{$\displaystyle 
b=c\cosh r+\eta c^\dagger\sinh r.
$}\]

Dies ist eine G-verträgliche kanonische Bogoliubov-Transformation. Sie
ergibt

\[\adjustbox{max width=\linewidth}{$\displaystyle 
\Delta'=\Delta\cosh2r+\kappa\sinh2r,\quad
\kappa'=\Delta\sinh2r+\kappa\cosh2r,
$}\] \[\adjustbox{max width=\linewidth}{$\displaystyle 
g'=g\cosh r+\lambda\sinh r,\quad
\lambda'=g\sinh r+\lambda\cosh r.
$}\]

Die additive Konstante ist 30 mal (Delta prime minus Delta). Bei g
Quadrat größer lambda Quadrat verschwinden kappa prime und lambda prime
gemeinsam genau dann, wenn

\[\adjustbox{max width=\linewidth}{$\displaystyle 
\mathcal I=\kappa(g^2+\lambda^2)-2\Delta g\lambda=0.
$}\]

Beispiel: Delta=1, kappa=4/5, g=2, lambda=1 und tanh r=-1/2 ergeben
Delta prime=3/5, g prime=Wurzel 3, kappa prime=lambda prime=0. Obwohl
die alte Zahl N nicht erhalten ist, gibt es dann die verborgene
kontinuierliche Erhaltung Nf plus zweimal Nc. Daher wäre auch die
Folgerung „N ist gebrochen, also ist das Fundament nur Z4`` ohne Prüfung
zu schnell. Für I ungleich null ist nur diese uniforme kanonische
Reduktion ausgeschlossen, nicht jede denkbare nichtlineare verborgene
Symmetrie.

Der sehr einfache reelle Quadraturansatz lambda=g bei kappa=0 ist
dagegen nicht durch eine endliche solche Transformation auf den nativen
Ansatz reduzierbar. Er ist ein prüfbarer Kandidat, kein hergeleitetes
Naturgesetz. Ein weiterer Herkunftstest ist besonders schlicht: Bei
freiem H0=epsilon Nf+Delta Nb haben die beiden kubischen Terme die
Frequenzen Delta minus 2 epsilon und Delta plus 2 epsilon. Am
verwendeten epsilon=0 gibt es keine schnelle gegen langsame Frequenz,
die das Weglassen des zweiten Kanals durch eine Rotating-Wave-Näherung
rechtfertigen würde. Ein exaktes U(1)-Prinzip könnte das dennoch
rechtfertigen; dieses Prinzip wäre dann auszuweisen und abzuleiten.

